\subsection*{The joint full-rigid and $3$-adic quotient still has a common image}

\statement{B.71}{Proved}{Both actual rigid kernels realize the full common Eichler principal group}
Let $R_T$ be B.64's subgroup of $H_T$, and let
\[
                       R_5=\langle\!\langle w_0\rangle\!\rangle_{H_5}
                                                               \tag{B.71.1}
\]
be the discrete normal closure of B.69's actual kernel word.
Use closures of these subgroups inside
$\widehat H_T,\widehat H_5$, respectively, embedded in
$\widehat P$ by B.28; no completion of the abstract
infinite-rank groups $R_T,R_5$ is being substituted.

Let $\mathscr R$ be the product of the twelve completed
pure seven-point target groups, and write
\[
 \Psi=(\widehat f_{\mathrm{rig}},q_\infty):\widehat P
          \longrightarrow\mathscr R\times
                         \operatorname{PSp}_{10}(\mathbb Z_3).
                                                               \tag{B.71.2}
\]
Here $q_\infty$, its principal identity fibre $\mathcal I$,
and the Eichler principal group $\mathcal U$ are as in B.64.
Then
\[
 \begin{split}
 \Psi(\overline{R_T})\cap\Psi(\overline{R_5})
       \cap(\{1\}\times\mathcal I)
       &=\{1\}\times\mathcal U,\\
 \Psi(\overline{H_T})\cap\Psi(\overline{H_5})
       \cap(\{1\}\times\mathcal I)
       &=\{1\}\times\mathcal U.
 \end{split}                                                    \tag{B.71.3}
\]
Thus the same nontrivial three-dimensional principal group
survives in the joint quotient containing all full rigid
targets and the entire $3$-adic homology action.
This is an equality in that joint quotient; it does not
give an equality of the source lifts in $\widehat P$.

\begin{proof}
By B.56 and the definition in B.64, $R_T$ lies in the
actual joint rigid kernel on $H_T$.
By B.69, $R_5$ lies in the actual joint rigid kernel
on $H_5$.  Both subgroups lie in their source commutator
subgroups.  Continuity therefore gives
\[
 \Psi(\overline{R_T})\subseteq\{1\}\times
                     \overline{q_\infty(R_T)},\qquad
 \Psi(\overline{R_5})\subseteq\{1\}\times
                     \overline{q_\infty(R_5)}.                \tag{B.71.4}
\]
The inclusions in (B.71.4) are equalities: each source
closure is compact, and its image is precisely the
closure of the corresponding discrete image.

It remains to strengthen B.64's cyclic lower bound
from $[H_5,H_5]$ to the actual rigid subgroup $R_5$.
Let $G_5$ be the connected derived algebraic centralizer
of $D$ on genus-five homology, and set
\[
 L_5=\operatorname{Lie}(G_5)\cap
                         \operatorname{Mat}_{10}(\mathbb Z_3),
 \qquad U_5=\exp(3L_5).
                                                               \tag{B.71.5}
\]
B.64 proves that $L_5$ is saturated of rank nine,
that this exponential is exactly the principal
subgroup of $G_5$, and that successive cubing generates
it from any basis of its first congruence quotient.
These assertions use the natural ten-dimensional
integral lattice and its $3L_5$ filtration.

The new actual words are particularly short to specify.
For
\[
 u\in\{1,f_0,f_1,f_2,f_3,f_5,f_0^2,f_1f_0,f_4f_2\},
 \qquad n_u=u\,w_0^{20}\,u^{-1},                            \tag{B.71.6}
\]
every $n_u$ belongs literally to $R_5\subseteq[H_5,H_5]$.
The exact program
\nolinkurl{next_3adic_symplectic/actual_cyclic_rigid_kernel.py}
recomputes all their homology matrices from the original
common-base eleven-strand braids.
It uses the integral branched-cover model of B.62,
checks its braid relations and symplectic inverses,
and makes no import from the discovery computation.
The projective reduction of $w_0$ modulo three has
order ten.  Squaring its tenth power ensures that
each $n_u$ has its actual linear lift congruent to $I$
modulo three, independently of a projective sign.
The nine matrices give residues
\[
                V_u=\frac{\rho(n_u)-I}{3}\pmod3
                       \in\operatorname{Mat}_{10}(\mathbb F_3)
                                                               \tag{B.71.7}
\]
of rank nine.
For an explicit rank witness, flatten matrix entries
row by row, numbered $0,\ldots,99$.
Rows $0,1,2,3,4,5,6,8,11$ of these nine columns form
an integer matrix with determinant $-52$, a unit
modulo three.
All nine actual source words, matrices, Lie residues
and the full minor are retained in
\nolinkurl{actual_cyclic_rigid_kernel_certificate.json},
with genuine standard output.
The reduced source word lengths are $200$, $202$ or $204$;
each has exactly zero integral source abelianization.

Every such commutator word maps into $G_5$, by B.64's
linear-centralizer and derived-group argument.
For a matrix in $I+3\operatorname{Mat}_{10}(\mathbb Z_3)$,
the logarithm divided by three reduces to (B.71.7).
Saturation identifies the reduction of $L_5$ with
its image in the ambient matrix lattice modulo three.
The nine independent residues are consequently a
basis of $L_5/3L_5$, not merely a basis of a smaller
unsaturated sublattice.

Apply B.64's explicit successive-cube argument to
these nine actual elements of $R_5$.
Their closed generated image contains every successive
congruence quotient of $U_5$, and hence
\[
                  U_5\subseteq\overline{q_\infty(R_5)}.
                                                               \tag{B.71.8}
\]
B.64 already proves the analogous inclusion
$U_{G_T}\subseteq\overline{q_\infty(R_T)}$.
The common Eichler principal group $\mathcal U$ is
contained in both $U_{G_T}$ and $U_5$.
Equation (B.71.4), with compact surjectivity, therefore
gives the lower inclusion in the first line of
(B.71.3), simultaneously in all full rigid coordinates
and at every $3$-power level.

Conversely, any common image of the two full source
closures has a matrix component centralizing $T$ and
$D$, hence their generated group $A_5$.
In $\mathcal I$ its matrix component belongs to
$\mathcal U$ by B.64's exact integral centralizer
calculation.  This gives the upper inclusion for
both lines of (B.71.3); the first lower inclusion
also proves the second one.

In particular no identification
$M_5=\overline{N_5}$ from B.70 is required.
The smaller explicitly generated subgroup $R_5$
already supplies the needed compact image, by closure
inside $\widehat H_5$.
The comparison is stronger than separate finite-image
tests because every chosen source lift kills every
full rigid coordinate exactly.
It still leaves the actual ambient intersection and
the residual centralizer $\mathcal K$ open.
\end{proof}
